% Generated draft: requires mathematical and source review.
\documentclass[11pt]{article}
\usepackage[margin=25mm]{geometry}
\usepackage{fontspec}
\setmainfont{Latin Modern Roman}
\usepackage{amsmath,amssymb,mathtools,unicode-math}
\setmathfont{Latin Modern Math}
\usepackage{enumitem,needspace,xurl,hyperref}
\setlength{\parindent}{0pt}
\setlength{\parskip}{5pt}
\setlength{\emergencystretch}{3em}
\newcommand{\N}{\mathbb{N}}
\newcommand{\Z}{\mathbb{Z}}
\newcommand{\Q}{\mathbb{Q}}
\newcommand{\R}{\mathbb{R}}
\newcommand{\C}{\mathbb{C}}
\newcommand{\F}{\mathbb{F}}
\newcommand{\A}{\mathbb{A}}
\newcommand{\PP}{\mathbb{P}}
\newcommand{\E}{\mathbb{E}}
\newcommand{\eps}{\varepsilon}
\newcommand{\dd}{\,\mathrm{d}}
\begin{document}
\section*{UnsolvedMath ID: 30003540}
\textbf{Source problem code:} catalogue entry
\subsection{Definitions and mathematical statement}
Let \(d\in\{1,2,3,\ldots\}\) and set \(p_d=2+4/d\). Use the normalization of the free Schrödinger equation
\[
 i\partial_tu+\Delta_xu=0,\qquad u(0,x)=f(x),\qquad (t,x)\in\R\times\R^d,
\]
where \(\Delta_x=\sum_{j=1}^d\partial_{x_j}^2\), \(f\in L^2(\R^d;\C)\), and \(u(t)=e^{it\Delta}f\). The source does not specify the equation's normalization; this convention fixes the meaning of the solution and its sharp constant. All norms below use Lebesgue measure, and the spacetime norm extends over all time and space:
\[
 \|u\|_{L^{p_d}(\R\times\R^d)}
 =\left(\int_{\R}\int_{\R^d}|u(t,x)|^{p_d}\,\dd x\,\dd t\right)^{1/p_d}.
\]
Define the sharp Strichartz constant by
\[
 C_d=\sup_{0\ne f\in L^2(\R^d;\C)}
 \frac{\|e^{it\Delta}f\|_{L^{p_d}(\R\times\R^d)}}{\|f\|_{L^2(\R^d)}}.
\]
An extremizer is a nonzero initial datum attaining this supremum, equivalently attaining equality in the inequality with constant \(C_d\).

\textbf{Question.} Do Gaussian initial data attain the sharp constant \(C_d\) for every positive integer \(d\)? Under the usual isotropic interpretation, a precise formulation is: with \(g_d(x)=e^{-|x|^2}\), does
\[
 \|e^{it\Delta}g_d\|_{L^{2+4/d}(\R\times\R^d)}
 =C_d\|g_d\|_{L^2(\R^d)}
 \qquad\text{hold for every }d\geq1?
\]
The source does not define its Gaussian class or specify whether its wording means existence of Gaussian extremizers or extremality of every member of a Gaussian class. The displayed formulation records the standard isotropic candidate; an interpretation encompassing arbitrary anisotropic or complex quadratic Gaussian profiles requires clarification. The question concerns attainment, and does not request a classification of all extremizers.

The supplied record labels the problem partially solved. Its dated assessment discusses hyperboloid endpoint constants and extremizers, including purported results for all non-endpoints in dimensions one and two, and says that the all-dimensional Gaussian-extremizer question remains open. The connection between those hyperboloid claims and this free Schrödinger inequality is not explained in the supplied evidence. These are supplied curation claims, not independently verified mathematical or literature conclusions.
\subsection{Short English statement}
For each positive integer dimension, consider the sharp spacetime Strichartz inequality for the free Schrödinger evolution of square-integrable initial data. The question asks whether Gaussian initial data attain its best constant in every dimension. The precise Gaussian class is unspecified; the displayed draft uses the standard isotropic Gaussian candidate. The supplied partial-solution label and hyperboloid literature discussion require independent review.
\subsection{Sources}
\begin{itemize}
\item Supplied local UnsolvedMath registry record: permanent ID \(30003540\), source problem code \texttt{catalogue entry}, title \emph{Gaussian Extremizers for Strichartz Inequalities}, collection \emph{Oberwolfach Reports Open Problems}. The record gives proposed year 2017 and supplies no named proposer.
\item Supplied source citation: \emph{Real Analysis, Harmonic Analysis, and Applications} (2018), Oberwolfach Reports, pp.\ 2101--2164, DOI \texttt{10.4171/owr/2017/34}. Original OWR report: \url{https://doi.org/10.4171/owr/2017/34}. The background describes these reports as records of workshops and their open-problem discussions. No problem-specific page is supplied.
\item Supplied literature reference, labeled ``Carneiro--Oliveira e Silva--Sousa, hyperboloid extremizers'': \url{https://arxiv.org/abs/1708.03826}. Its precise relationship to the stated Schrödinger Gaussian-extremizer question is not established by the record.
\item Literature assessment dated exactly \texttt{2026-08-21}: supplied dated curation, not independently verified. The record labels the status partially solved and states that hyperboloid endpoint constants and extremizers are known in several low-dimensional cases, including all non-endpoints in dimensions one and two, while the all-dimensional Gaussian-extremizer question remains open. The background explicitly requires independent verification of time-sensitive status and cited claims.
\item The registry supplies statement status \texttt{corrected\_verified}, accompanied by a claim of manual semantic, typographic, and LaTeX review. This is a supplied review label, not independent verification of the mathematical formulation or literature assessment.
\end{itemize}
\end{document}
